% from ega4-17-fr.tex
\oldpage[IV-4]{87}

  \item[c)] Les images des $d_{X/S}g_i$ dans $(\Omega^1_{X/S})_x\otimes_{\sh{O}_{X,x}}k(x)$ sont des éléments linéairement indépendants sur $k(x)$, et les $(g_i)_x$ engendrent $\sh{J}_x$.
  \item[d)] Il existe un voisinage ouvert $U$ de $x$ dans $X$ et des sections $g_j$ ($r+1\leq j\leq n$) de $\sh{O}_X$ au-dessus de $U$ telles que le $S$-morphisme $u:U\to S[T_1,\ldots,T_n]$ correspondant à l'homomorphisme $\sh{O}_S^n\to f_*(\sh{O}_U)$ défini par les sections $g_i|U$ ($1\leq i\leq r$) et $g_j$ ($r+1\leq j\leq n$) (II, 1.2.7) soit un morphisme étale, et que l'image réciproque du sous-préschéma $Y'=S[T_{r+1},\ldots,T_n]$ par ce morphisme soit le préschéma induit par $j(Y)$ sur l'ouvert $U\cap j(Y)$.
\end{enumerate}
\end{corollary}

Comme $(g_i')_x$ est l'image canonique de $(g_i)_x$, l'équivalence de a) et b) résulte du lemme de Nakayama, $\sh{J}_x$ étant de type fini et $\sh{J}_x/\sh{J}_x^2$ un $(\sh{O}_{X,x}/\sh{J}_x)$-module libre (17.10.4) (Bourbaki, \emph{Alg. comm.}, chap. II, § 3, n\textsuperscript{o} 2, prop. 5). En vertu de (17.12.1, b)), $\sh{J}_x/\sh{J}_x^2$ s'identifie canoniquement à un facteur direct du $(\sh{O}_{X,x}/\sh{J}_x)$-module libre de rang $n$, $(\Omega^1_{X/S})_x\otimes_{\sh{O}_{X,x}}(\sh{O}_{X,x}/\sh{J}_x)$, et l'équivalence de b) et c) résulte de Bourbaki, \emph{loc. cit.} En outre, si a) est vérifiée, $(g_i')_x$ est ainsi identifié à $(d_{X/S}g_i)_x\otimes1$ pour $1\leq i\leq r$; le $\sh{O}_{X,x}$-module $(\Omega^1_{X/S})_x$ étant libre de rang $n$, il résulte encore de Bourbaki, \emph{loc. cit.}, qu'en remplaçant au besoin $X$ par un voisinage ouvert de $x$, on peut supposer qu'il existe $n-r$ sections $g_i$ ($r+1\leq i\leq n$) de $\sh{O}_X$ au-dessus de $X$, telles que les $(d_{X/S}g_i)_x$ pour $1\leq i\leq n$ forment une base de $(\Omega^1_{X/S})_x$; le fait que le morphisme correspondant $u$ soit étale au point $x$ résulte alors de (17.11.3); pour la même raison, le morphisme $u^{-1}(Y')\to Y'$, restriction de $u$, est étale au point $x$; en remplaçant $X$ par un voisinage de $x$, on peut supposer ces deux morphismes étales. En outre, il est immédiat que $Y$ (identifié, pour simplifier, au sous-préschéma fermé $j(Y)$ de $X$) est un sous-préschéma fermé de $u^{-1}(Y')$, et en vertu du choix des $g_i$ pour $i>r$ et de (17.2.5) et (17.11.2), la restriction à $Y$ de $u$ peut encore être supposée étale. On en déduit que pour tout $y'\in Y'$, l'immersion $Y\cap u^{-1}(y')\to u^{-1}(y')$ est ouverte, d'où l'on conclut à l'aide de (17.9.6) que $Y\to u^{-1}(Y')$ est une immersion ouverte; cela prouve que a) entraîne d). La réciproque découle aussitôt de (17.11.3).

\begin{corollary}[17.12.3]
\label{IV.17.12.3}
Soient $f:X\to S$ un morphisme localement de présentation finie, $u:S\to X$ une $S$-section de $X$. Pour que $f$ soit lisse en un point $x\in X$, il faut et il suffit que $u$ soit une immersion quasi-régulière dans un voisinage de $s=f(x)$.
\end{corollary}

La conclusion résulte de (17.12.1), puisque $f\circ u=1_S$ est lisse.

\begin{proposition}[17.12.4]
\label{IV.17.12.4}
Tout morphisme lisse $f:X\to S$ est différentiellement lisse (autrement dit (16.10.5) $\Delta_f:X\to X\times_SX$ est une immersion quasi-régulière). \emph{En particulier, les $\sh{P}_f^n=\sh{P}_{X/S}^n$ et les $\operatorname{gr}_n(\sh{P}_{X/S})$ sont des $\sh{O}_X$-Modules localement libres de type fini.}
\end{proposition}

Il suffit de remarquer que le morphisme structural $p_1:X\times_SX\to X$ est lisse (17.3.3, (iii)) et que $\Delta_f$ est une $X$-section pour $p_1$; la conclusion résulte de (17.12.3).

\begin{proposition}[17.12.5]
\label{IV.17.12.5}
Soit $f:X\to Y$ un morphisme localement de présentation finie. Les conditions suivantes sont équivalentes:
\begin{enumerate}
  \item[a)] $f$ est différentiellement lisse.
  \item[b)] Pour tout morphisme $g:Y'\to Y$, si l'on pose $X'=X\times_Y Y'$ et $f'=f_{(Y')}:X'\to Y'$, alors, pour toute $Y'$-section $s'$ de $X'$, $f'$ est lisse en tous les points de $s'(Y')$.
  \item[c)] La seconde projection $p_2:X\times_YX\to X$ est lisse en tous les points de la diagonale $\Delta_f(X)$.
\end{enumerate}
\end{proposition}

